% ----------------------------------------------------------
\chapter{Scripts de Processamento e Análise de Dados}
\label{ap:scripts_processamento}
% ----------------------------------------------------------

Este apêndice descreve a estrutura de scripts desenvolvida para a ingestão, limpeza, padronização, análise estatística e classificação de gargalos da base de dados telemétrica das redes Ruckus e UniFi. 

Todo o código-fonte correspondente e as instruções detalhadas de configuração do ambiente de execução estão disponíveis publicamente no repositório do projeto hospedado no GitHub, acessível por meio do seguinte endereço:
\begin{center}
    \url{https://github.com/locdown2311/DECSI-COSI-Modelo-Monografia/tree/v2/scripts}
\end{center}

Abaixo são detalhados os principais scripts do pipeline executados para a consolidação dos resultados apresentados nesta monografia:

\begin{enumerate}
    \item \textbf{\texttt{unify\_logs.py}}: Consolida os arquivos de log diários gerados pelas controladoras em um único arquivo bruto unificado para cada fabricante, aplicando filtros iniciais de integridade de data.
    \item \textbf{\texttt{dataset\_validator.py}}: Realiza o diagnóstico inicial de qualidade de dados brutos consolidados, identificando falhas de amostragem temporal, taxas de medições incompletas e preenchendo ou descartando registros corrompidos.
    \item \textbf{\texttt{data\_standardizer.py}}: Realiza a limpeza fina de registros inválidos (como MACs vazios ou instâncias sem leitura), converte as grandezas para escalas físicas comuns (e.g., potência de sinal recebido em dBm e vazão em Mbps) e mapeia os nomes legíveis dos Access Points cadastrados.
    \item \textbf{\texttt{descriptive\_analyzer.py}}: Executa o cálculo das estatísticas descritivas básicas (médias, medianas, quartis e desvios padrão) de todos os parâmetros físicos e plota histogramas e diagramas de caixa (\textit{boxplots}) comparativos.
    \item \textbf{\texttt{visualizador\_dados.py}}: Gera gráficos temporais de distribuição de clientes únicos ao longo dos dias e os mapas de calor de densidade horária de dispositivos associados por fabricante.
    \item \textbf{\texttt{analisador\_correlacoes.py}}: Computa as matrizes de correlação cruzada usando os coeficientes de Pearson e Spearman, gerando os mapas de calor gerais e gráficos de dispersão específicos entre variáveis críticas.
    \item \textbf{\texttt{analisador\_testes.py}}: Aplica os testes de normalidade de Shapiro-Wilk e de significância estatística U de Mann-Whitney para validação das hipóteses, acompanhado do cálculo do tamanho de efeito operacional (\textit{effect size} $r$ e $d$).
    \item \textbf{\texttt{classificador\_gargalos.py}}: Executa a classificação heurística de gargalos operacionais de rede (G1 a G5), realiza o agrupamento de Access Points via algoritmo K-Means e gera os gráficos correspondentes de distribuição de carga.
    \item \textbf{\texttt{run\_pipeline.py}}: Script orquestrador central que executa todo o pipeline de processamento na sequência lógica correta, gerando as tabelas em CSV e convertendo os relatórios finais.
    \item \textbf{\texttt{convert\_md\_to\_tex.py}}: Script utilitário responsável por ler os relatórios em Markdown, aplicar as regras de escape LaTeX e formatação brasileira de números (decimais por vírgula, milhares por ponto) e exportar os arquivos textuais finais.
\end{enumerate}
